% =====================================================================
%  BAB 2 — TINJAUAN PUSTAKA
% =====================================================================

\chapter{Tinjauan Pustaka}

% ── 2.1 Landasan Teori ───────────────────────────────────────────────
\section{Landasan Teori}
\label{sec:2.1}

\subsection{[Konsep Utama 1]}
\label{sec:2.1.1}

[Penjelasan konsep utama pertama yang mendasari penelitian ini. Definisi, cara kerja,
dan relevansinya terhadap penelitian. Sertakan rujukan ke literatur primer.]

\subsection{[Konsep Utama 2]}
\label{sec:2.1.2}

[Penjelasan konsep utama kedua.]

\subsection{[Konsep Utama 3]}
\label{sec:2.1.3}

[Penjelasan konsep utama ketiga.]

% Contoh persamaan matematis
% Misalnya: F1-score
\begin{equation}
  F_1 = 2 \cdot \frac{\text{Presisi} \times \text{Recall}}{\text{Presisi} + \text{Recall}}
  \label{eq:f1}
\end{equation}

% Contoh tabel perbandingan
% \begin{table}[H]
%   \centering
%   \caption{Perbandingan [hal yang dibandingkan]}
%   \label{tab:2.perbandingan}
%   \begin{tabular}{llll}
%     \toprule
%     \textbf{Metode} & \textbf{Kelebihan} & \textbf{Kekurangan} & \textbf{Sumber} \\
%     \midrule
%     {[Metode A]} & [+] & [-] & [Penulis, Tahun] \\
%     {[Metode B]} & [+] & [-] & [Penulis, Tahun] \\
%     \bottomrule
%   \end{tabular}
% \end{table}

% ── 2.2 Penelitian Terdahulu ─────────────────────────────────────────
\section{Penelitian Terdahulu}
\label{sec:2.2}

[Ulas penelitian-penelitian yang relevan secara kritis. Identifikasi kontribusi, metode,
dataset, dan hasil dari setiap penelitian. Gunakan paragraf tematis, bukan hanya daftar
rangkuman.

Struktur yang baik: kelompokkan penelitian berdasarkan pendekatan, bukan kronologi.
Tunjukkan bagaimana setiap penelitian berkaitan dengan gap yang Anda isi.]

\subsection{[Kelompok Penelitian 1: Pendekatan X]}
\label{sec:2.2.1}

[Ulas beberapa penelitian yang menggunakan pendekatan X. Sintesis persamaan dan
perbedaannya.]

\subsection{[Kelompok Penelitian 2: Pendekatan Y]}
\label{sec:2.2.2}

[Ulas pendekatan Y dan kelemahannya yang menjadi motivasi penelitian Anda.]

% ── 2.3 Posisi Penelitian ────────────────────────────────────────────
\section{Posisi Penelitian}
\label{sec:2.3}

Berdasarkan tinjauan terhadap penelitian-penelitian yang telah diuraikan, Tabel~\ref{tab:2.posisi}
merangkum posisi penelitian ini dibandingkan karya terdahulu.

\begin{table}[H]
  \centering
  \caption{Posisi penelitian ini terhadap penelitian terdahulu}
  \label{tab:2.posisi}
  \begin{tabular}{p{2.9cm} p{1.8cm} p{1.8cm} p{1.8cm} p{1.8cm} p{1.8cm}}
    \toprule
    \textbf{Penelitian} & \textbf{[Aspek 1]} & \textbf{[Aspek 2]} & \textbf{[Aspek 3]}
                        & \textbf{[Aspek 4]} & \textbf{[Aspek 5]} \\
    \midrule
    {[Peneliti A, Tahun]} & \checkmark & -- & \checkmark & -- & -- \\
    {[Peneliti B, Tahun]} & \checkmark & \checkmark & -- & -- & -- \\
    \textbf{Penelitian ini} & \checkmark & \checkmark & \checkmark & \checkmark & \checkmark \\
    \bottomrule
  \end{tabular}
\end{table}

[Paragraf yang menjelaskan kebaruan (\emph{novelty}) penelitian ini berdasarkan tabel di atas.
Apa yang belum dilakukan orang lain dan yang Anda lakukan?]
